\chapter*{Streszczenie}
Prezentowana praca dyplomowa skupia się na przedstawieniu różnic w użytkowaniu dużych modeli językowych, będących częścią wybranych konfiguracji systemów RAG. Autor zaprezentował przebieg i wyniki badania, które miało na celu sprawdzić wydajność tych algorytmów podczas przetwarzania pytań, dotyczących zewnętrznych nieustrukturyzowanych i ustrukturyzowanych źródeł danych. W tym celu, utworzono oprogramowanie do testów i interfejs, umożliwiający prowadzenie konwersacji z dwoma wybranymi optymalnymi systemami.

Część teoretyczna składa się z wprowadzenia do przetwarzania języka naturalnego, dużych modeli językowych oraz techniki generowania wspomaganego wyszukiwaniem. Scharakteryzowano między innymi: zarys historyczny chatbotów, sposób w jaki maszyny rozumieją ludzki język oraz jak działają LLMy i dlaczego są one chętnie wykorzystywane w ramach RAG.

Część związana z przeglądem literatury opisuje w jaki sposób naukowcy podchodzą do ewaluacji systemów RAG. Pierwsza przywołana publikacja zaprezentowała skalowalny wielodokumentowy test wydajnościowy MEBench. W drugim artykule autorzy pochylili się nad benchmarkiem RGB, w ramach którego przetestowano systemy RAG tysiącem pytań z czterech kategorii: odporność na szum, odmowa odpowiedzi, integracja informacji i odporność na błędne fakty. Natomiast, w trzeciej publikacji opisano benchmark NL-to-SQL, składający się łącznie z 360 łatwych, umiarkowanych i trudnych pytań odnoszących się do bazy danych SQL.

W części badawczej przedstawiono autorskie oprogramowanie umożliwiające automatyzację testowania trzydziestu różnych systemów RAG. Autor dokładnie opisał jak powstały dane nieustrukturyzowane w postaci pliku PDF i tekstowego  oraz dane ustrukturyzowane w postaci bazy danych SQL i pliku CSV. Inspirując się przestudiowaną literaturą, zaproponowano dwa zestawy ośmiu pytań, dedykowanych dwóm typom wyżej wymienionych źródeł informacji.

Część, w której skupiono się na analizie wyników badania, prezentuje porównanie systemów RAG pod kątem trafności generowanych odpowiedzi, kosztu ich uzyskania, czasu generowania oraz przepustowości tokenów. W ostatnim kroku, wybrano dwa najwydajniejsze systemy, po jednym z każdej kategorii. Następnie zintegrowano je z graficznym interfejsem użytkownika, co umożliwiło prowadzenie z nimi konwersacji.



\bigskip
\bigskip

\noindent \textbf{Słowa kluczowe}: \\
NLP, Chatbot, sztuczna inteligencja, LLM, RAG, agent AI, dane nieustrukturyzowane, dane ustrukturyzowane, PDF, plik tekstowy, SQL, CSV, automatyzacja, GUI, scraping, embedder, wektorowa baza danych


\bigskip

\noindent \textbf{Dziedzina nauki i techniki, zgodnie z wymogami OECD}: \\
Nauki inżynieryjne i techniczne, informatyka techniczna i telekomunikacja